% !TEX root = ../Dissertation.tex

\chapter{{Schlussbemerkung und Ausblick}} \label{Gesamtfazit}
Den in diesem Buch vorgestellten empirischen Untersuchungen liegt die Beobachtung zugrunde, dass es im Deutschen eine Vielzahl von Intensitätspartikeln gibt, die sich in mehreren Gesichtspunkten zu unterscheiden scheinen. Dabei sind ältere Intensivierer wie \textit{arg} oder \textit{voll} naturgemäß nicht nur lexikalisierter als die neueren, sondern dem subjektiven Anschein nach auch weniger expressiv. Neuere Intensivierer wie \textit{krass} oder \textit{mega} zeichnen sich neben ihrer skalierenden Funktion auch durch eine besondere pragmatische Schlagkraft, d.\,h. eine überdurchschnittlich ausgeprägte Expressivität, aus, mit der Sprecher ihre persönliche Einstellung zum bezeichneten Sachverhalt zum Ausdruck bringen können. Schwerpunkt der Untersuchungen waren einige in der theoretischen Literatur vertretene Thesen, die zum ersten Mal kritisch hinterfragt und auf ihre Gültigkeit hin überprüft wurden. So wird der Variantenreichtum von Intensivierern bspw. durch starke Erneuerungstendenzen erklärt, die durch die mit den Ausdrücken einhergehende Affektivität bedingt seien \mbox{\citep[vgl.][650]{Bussmann.2002}}. Diese nutzt sich mit frequentem und persistentem Gebrauch sukzessive ab, so die Aussagen von \mbox{\citet[2]{Hauschild.1899}} und \mbox{\citet[59]{Tobler.1868}}. Ausgehend von diesen Annahmen lässt sich folgende Forschungshypothese ableiten: Je frequenter oder länger ein Intensivierer im Gebrauch ist, je vorhersehbarer er also ist, desto niedriger ist sein Expressivitätsgrad (vgl. H1-Freq und \mbox{H1-Temp}). Nach \mbox{\citet[274]{Baumgarten.1908}} resultiert die Abnutzung in einem ständigen Bedarf an neuen, unverbrauchten Ausdrücken, mit denen Sprecher den Ausfall altgedienter Formen kompensieren und expressiv bleiben können (\mbox{vgl. H2}). Im Zentrum des Empirieteils stand die Überlegung, dass sich der Expressivitätsgrad von Intensivierern skalar erfassen lässt, sodass ihnen in Abhängigkeit der mit dem jeweiligen Ausdruck verbundenen Intensität auf der Skala ein bestimmter Wert zugeordnet werden kann. Diese Idee wurde u.\,a. auf folgende Hypothese übertragen: Wird ein Ausdruck als stark expressiv wahrgenommen, so wird er in Bezug auf Skalarität (vgl. H3) und Sprechereinstellung (vgl. H4) höher verortet als ein wenig expressiver Ausdruck, wie \citet[22]{Paradis.1997} und \mbox{\citet[133]{Gutzmann.2019a}} sowie \mbox{\citet[315]{Lang.1983}} und \mbox{\citet[57]{Ortner.2014}} postulieren. Die durch expressive Intensivierer kommunizierte Einstellungsspiegelung ist nach Auffassung von \citet[135]{Gutzmann.2019a} evaluativ und äußert sich in einer positiven oder negativen Bewertung der im Satz beschriebenen Sachlage. Dagegen drücken deskriptive Intensivierer keine Sprechereinstellung aus und werden dementsprechend als überwiegend neutral empfunden. Dies warf die Frage auf, wie die Einstellungspolarität beurteilt wird, wenn das modifizierte Element ein bedeutungsloses Pseudoadjektiv wie \textit{fongig} ist. In diesem Zusammenhang wurde im Vorfeld der Untersuchung erwogen, dass die Einschätzung der Einstellungspolarität nicht durch das modifizierte Pseudoadjektiv, sondern durch die evaluative Polarität des Ursprungslexems gesteuert wird: Ist ein Intensivierer aufgrund semantischer Transparenz negativ konnotiert (bspw. \textit{furchtbar} oder \textit{schrecklich}), so wird die Sprechereinstellung zugleich verstärkt als negativ wahrgenommen. Gleiches gilt für positiv konnotierte Lexeme wie \textit{mega} oder \textit{super}, die die Einstellungspolarität tendenziell in die positive Richtung lenken (vgl. H5). Außerdem wurden in dieser Studie die in der Literatur vertretenen intraklassischen Unterschiede untersucht, die sich womöglich in Bezug auf Skalarität und Sprechereinstellung zeigen (vgl. z.\,B. \citealt[87]{Biedermann.1969} und \mbox{\citealt[408]{Breindl.2007}}). Diesbezüglich wird angenommen, dass sich die Mitglieder einer Intensiviererklasse in Abhängigkeit des durch sie ausgedrückten Intensitätsgrades auf einer kontinuierlichen Skala anordnen lassen \mbox{wird (vgl. H6)}. Nach \mbox{\citet[133]{Gutzmann.2019a}} ist zudem davon auszugehen, dass die Expressivität der Intensivierer Einfluss auf deren Skalaritätsgrad nimmt, der in der Folge als stärker wahrgenommen wird. \par
Die skizzierten Hypothesen werden in der Literatur als gegeben, da intuitiv einleuchtend, betrachtet, gleichwohl deren Gültigkeit bislang nicht empirisch abgesichert ist. Dieses Ziel wurde durch eine breit angelegte Untersuchung und eine systematisch aufeinander aufbauende Kombination aus verschiedenen linguistischen Standardmethoden verfolgt, durch die erstmals gesicherte Aussagen über das Zutreffen der Annahmen gemacht werden können. Durch die empirische Herangehensweise schließt dieses Buch die Lücke, die auf diesem Gebiet vorherrscht, sind die bisherigen Forschungsbeiträge von wenigen Ausnahmen abgesehen doch theoretischer Natur.


Die Ergebnisse der korpuslinguistischen und experimentellen Untersuchungen belegen die Gültigkeit beinahe aller zur Diskussion stehenden Hypothesen. Die Hypothese H1-Freq nach \citet[2]{Hauschild.1899}, nach der inflationärer Gebrauch von Intensivierern zulasten ihrer Expressivität geht, wird durch die Ergebnisse des in Kapitel \ref{Experiment0} geschilderten \mbox{Experiments 1} gestützt. In diesem findet sich ein signifikanter Unterschied zwischen den im Korpus neu aufgetretenen und den bereits länger bestehenden Intensivierern. Dies lässt darauf schließen, dass eine hohe Gebrauchsfrequenz zu einem bedeutend niedrigeren Grad an Sprechereinstellung führt, sie also tatsächlich abträglich für die Expressivität von Intensivierern ist. Das Auftreten der neuen Intensivierer \textit{hammer}, \textit{krass} und \textit{mega} im Korpus bestätigt weiterhin die Hypothese H2 nach \citet[274]{Baumgarten.1908}, nach der ständig neue Ausdrücke erforderlich sind, die den Ausfall verbrauchter Intensivierer kompensieren. Die Ergebnisse von Experiment 1 bekräftigen außerdem die Gültigkeit der Hypothesen H3 und H4 nach \citet[22]{Paradis.1997} und \citet[133]{Gutzmann.2019a} sowie \citet[315]{Lang.1983} und \citet[57]{Ortner.2014}. So wurden Sätze mit expressivem Intensivierer wie \textit{mega} oder \textit{super} von den Versuchspersonen im Hinblick auf Skalarität und Sprechereinstellung auf der Skala signifikant höher positioniert als Sätze, die den deskriptiven Intensivierer \textit{sehr} enthielten oder im Experiment ohne Intensivierer präsentiert wurden. In Bezug auf die Hypothese H5 nach \citet[135]{Gutzmann.2019a} lässt sich festhalten, dass vor allem diejenigen Sätze als evaluativ beurteilt werden, die mit einem hohen Grad an Skalarität oder Sprechereinstellung assoziiert werden. Zudem legen die Ergebnisse nahe, dass die Polarität des Ursprungslexems offenbar auf die Interpretation der Sprechereinstellung einwirkt. Diese Annahme beruht auf der Beobachtung, dass Intensivierer mit klar negativem Ursprung wie \textit{furchtbar} oder \textit{schrecklich} von den Versuchspersonen verstärkt als negativ beurteilt wurden, wohingegen die Interpretation positiv konnotierter Intensivierer wie \textit{mega} oder \textit{super} vornehmlich in die positive Richtung tendierte. Als neutral wurde abgesehen vom opaken deskriptiven Intensivierer \textit{sehr} bloß die unmarkierte Variante des Satzes eingeschätzt, die als Kontrollstufen in die Untersuchung miteinbezogen wurden. Darüber hinaus liefern die in Kapitel \ref{Experimente} vorgestellten Experimente 2a und 2b Evidenz für die Korrektheit der Hypothese H6 im Sinne von \citet[87]{Biedermann.1969} und \citet[408]{Breindl.2007}, die von klasseninternen Unterschieden der durch die Ausdrücke vermittelten Intensität ausgeht. Dabei treten innerhalb beider Intensiviererklassen sowohl bei Skalarität als auch bei Sprechereinstellung die vermuteten graduellen Abstufungen auf, deren Ursache semantisch begründet zu sein scheint.


%Dies deutet auf eine natürliche Realität des Unterschiedes zwischen deskriptiven und expressiven Intensivierern hin, der keineswegs zufallsbedingt ist.


Durch die Untersuchungsergebnisse nicht gestützt wird hingegen die Hypothese \mbox{H1-Temp} nach \mbox{\citet[59]{Tobler.1868}}, nach der langer Gebrauch einen negativen Effekt auf die Expressivität von Ausdrücken hat. Wie in \mbox{Abschnitt \ref{abs:Diskussion22}} erläutert, ist dieser Befund jedoch mit Vorsicht zu sehen, ist die Gruppe der neuen Intensivierer mit lediglich drei Mitgliedern (d.\,h. \textit{hammer}, \textit{krass} und \textit{mega}) doch stark unterrepräsentiert, weshalb die Resultate an dieser Stelle als nur wenig aufschlussreich gewertet werden können. Trotz der vorliegenden Ergebnisse ist es keineswegs abwegig anzunehmen, dass neben zu häufigen auch lange verwendete Intensivierer weniger überraschend und daher weniger expressiv sind als neuere. Um stichhaltige Aussagen über die Gültigkeit von \mbox{H1-Temp} zu machen, empfehlen sich weiterführende Experimente, die ein ausgewogenes Verhältnis aus alten und neuen Ausdrücken berücksichtigen. Des Weiteren kann auch die fehlende Korrelation bei der in \mbox{Kapitel \ref{Experimente}} dargelegten Untersuchung bloß mit Vorsicht interpretiert werden. Auch wenn kein korrelativer Zusammenhang zwischen der in \mbox{Experiment 2a} erfragten Skalarität und der in Experiment 2b ermittelten Sprechereinstellung nachgewiesen werden kann, bedeutet dies nicht, dass die beiden Merkmale grundsätzlich nicht korrelieren. Auch hier bedarf es einer größeren Datenmenge als derjenigen, die in die Korrelationsanalyse einbezogen wurde.


Neben den gewonnenen Ergebnissen gibt es noch einige Fragen, die über die Untersuchung hinaus offen sind und als Gegenstand für Folgestudien dienen können. Diese werden im Anschluss angeschnitten. Obendrein werden bestehende Forschungsdesiderate sowie abschließende Gedanken und Ideen aufgezeigt, die ebenfalls relevant für zukünftige experimentelle Untersuchungen sein können. Dass auf dem Gebiet der Expressivität vor allem letztere nur spärlich vertreten sind, wurde mehrfach hervorgehoben.


Ein noch nicht abschließend geklärter Punkt bezieht sich auf die in Abschnitt \ref{Intensivierertypen} angesprochene Möglichkeit zur DP-externen Positionierung von expressiven Intensivierern (z.\,B. \textit{total} oder \textit{voll}), die u.\,a. \mbox{\citet[]{Gutzmann.2012}} thematisieren. Vor diesem Hintergrund habe ich beanstandet, dass die Autor:innen besagte Eigenart uneingeschränkt auf die gesamte Klasse der expressiven Intensivierer beziehen. Wie \mbox{Beispiel \ref{Haha}} allerdings demonstrierte, ist die Behauptung in dieser Form deutlich zu undifferenziert. So scheint es graduelle Abstufungen hinsichtlich der Akzeptabilität zu geben, die \citet[]{Gutzmann.2012} komplett außer Acht lassen. Aus dieser intuitiven Einschätzung heraus ergibt sich die Notwendigkeit, das Phänomen der DP-externen Positionierung mittels experimenteller Methoden wie Akzeptabilitätsurteilen in Form von Skalenratings im Hinblick auf den Grad an Natürlichkeit hin zu untersuchen. Dabei könnten auch die äußeren Rahmenbedingungen hinzugezogen und regionale bzw. diatopische, soziolektale sowie prosodische Präferenzen berücksichtigt werden, die gegebenenfalls Einfluss auf die Beurteilung nehmen. Genauso bietet die in Experiment 1 überprüfte These von \mbox{\citet[135]{Gutzmann.2019a}} genügend Potenzial für weitere Untersuchungen. Nach dieser ist die Einstellungsbekundung, die expressive Intensivierer kommunizieren, evaluativen Charakters, weswegen sie sich in Form einer positiven oder negativen Bewertung des beschriebenen Sachverhalts niederschlägt. Weil sich die in \mbox{Kapitel \ref{Experiment0}} präsentierte Untersuchung ausschließlich auf die Modifikation semantisch leerer Pseudoadjektive wie \textit{fongig} und \textit{kieflich} bezog, ist Gutzmanns Bemerkung im Zusammenhang mit realen Wörtern noch immer unerforscht. Somit könnte auch sie Gegenstand künftiger experimenteller Studien mit einem ähnlichen Untersuchungsaufbau wie dem in diesem Buch beschriebenen sein, liefert eine komparative Konstant-Summen-Skala doch feinkörnige Abstufungen zwischen den Auswahloptionen, die weitere Rückschlüsse auf die Rolle der Semantik und demnach das Zutreffen dieser \mbox{Aussage zulassen}.


Daneben konnte auch einigen im Vorfeld der Untersuchungen ausgearbeiteten Ideen nicht nachgegangen werden, die sich jedoch ebenso als fruchtbar für weitere empirische Überprüfungen erweisen können. Die erste betrifft den Faktor des Alters, der schon in einigen Korpusstudien Beachtung findet: Je mehr ältere Personen einen Intensivierer kennen, desto länger gibt es den Ausdruck vermutlich in ebendieser Funktion. Dies beruht auf der Auffassung, dass sich neue Ausdrücke erst einmal im Sprachgebrauch jüngerer Menschen einlagern, die dahingehend offener und experimentierfreudiger sind als ältere Menschen (vgl. u.\,a. \citealt[]{Androutsopoulos.1998,Androutsopoulos.1999}; \citealt[]{Paradis.2000a}; \citealt[]{Lorenz.2002}; \citealt[]{Kirschbaum.2002a}; \citealt[]{Tagliamonte.2008}; \citealt[]{Tagliamonte.2014a}; \citealt[]{Neuland.2018}; \citealt[]{Stratton.2020a}; \citealt[]{Tagliamonte.2020}). Demzufolge ist anzunehmen, dass es altersbedingte Präferenzen und Unterschiede gibt, die weitere Anhaltspunkte für den diachronen Gebrauch und die Abnutzung von Intensivierern bieten können. Außerdem wird in der theoretischen Literatur seit Jahren die Annahme vertreten, dass Jüngere grundsätzlich mehr Intensivierer als Ältere verwenden (bspw. \mbox{\citealt[]{Paradis.2000a}}; \mbox{\citealt[]{Ito.2003}}). Ob diese These der Realität entspricht oder ob letztere lediglich auf andere, (heute) möglicherweise weniger ausdrucksstarke Formen zurückgreifen, bleibt zukünftig zu klären. Für die Untersuchung würden sich z.\,B. Experimente mit Akzeptabilitätsurteilen anbieten, in denen eine Personenstichprobe die Natürlichkeit oder Geläufigkeit verschiedener Intensivierer skalenbasiert beurteilt. In Anbetracht des Alters argumentiert \mbox{\citet[229]{Pittner.1989}} ferner, dass sich die mit einem Ausdruck verbundene Konnotation in Abhängigkeit von Sprecher und Äußerungskontext verändern könne: Während \textit{sau} für Erwachsene stärker negativ konnotiert sei, werde es von Jugendlichen vielmehr mit etwas Positivem in Verbindung gebracht (\textit{sauteuer}, \textit{saudumm} vs. \textit{saugeil}, \textit{saustark}) \citep[vgl.][229]{Pittner.1989}. Auch diese Einschätzung bedarf zweifellos empirischer Überprüfung, die, soweit mir bekannt, bislang nicht in Angriff genommen wurde. Überdies ist vorstellbar, dass der Altersfaktor ebenfalls mit kognitiven Verarbeitungsunterschieden einhergeht: Je neuer, überraschender und unbekannter ein Intensivierer ist, desto kostenträchtiger ist mutmaßlich seine Verarbeitung. Daraus leitet sich die Überlegung ab, dass sich die potenziellen alters- und personenspezifischen Verarbeitungsunterschiede mithilfe von neurophysiologischen, zeitsensitiven oder behavioralen Verfahren wie EEG, Self-Paced Reading oder Eye-Tracking erfassen lassen. Hier könnten Unterschiede im EKP, in Lesezeiten oder Pupillenreaktionen ein Indikator für eine erhöhte Rechenleistung und einen höheren kognitiven Arbeitsaufwand sein. Ebenso wie das Alter könnten auch Geschlecht und sozioökonomischer Status den Gebrauch von Intensivierern mitbestimmen. So gibt es bereits einige Korpusuntersuchungen, deren Ergebnisse geschlechtsspezifische und statusabhängige Präferenzen nahelegen (u.\,a. \citealt[]{Macaulay.1995}; \citealt[]{Ito.2003}; \citealt[]{Tagliamonte.2005}; \citealt[]{Tagliamonte.2008}; \citealt[]{Setayesh.2018}; \citealt[]{Stratton.2020b}; \citealt[]{Stratton.2020a}; \citealt[]{Stratton.2022}; \citealt[]{Stratton.2022b}). Diese Fragestellungen sind meines Wissens nach gegenwärtigem Forschungsstand noch nicht erschöpfend beantwortet und sollten unbedingt in gezielte experimentelle Untersuchungen überführt werden. \par 
Eine völlig andere Richtung, die im Rahmen der Untersuchung ebenfalls nicht verfolgt werden konnte, knüpft an die Etymologie von Intensitätspartikeln an. So entsteht der Eindruck, dass moderne Intensivierer häufig aus anderen Sprachen abgeleitet sind, was z.\,B. an \textit{krass} oder \textit{mega} zu erkennen ist. Dagegen sind ältere Intensivierer wie \textit{furchtbar} oder \textit{voll} meist nativ und deutschen Ursprungs. Somit steht die Frage im Raum, welche Faktoren die fremdsprachlichen Tendenzen bei der Intensivierung konditionieren. In diesem Kontext ist denkbar, dass sich unsere digitale, computerzentrierte und globalisierte Zeit auf die Entstehung von Intensivierern auswirkt, wie bspw. \mbox{\citet[354]{Claudi.2006}} als Erklärung für die Beliebtheit von \textit{giga} und \textit{mega} anführt. Darüber hinaus zeichnet sich über die Zeit ein morphologischer Wandel ab, der sich durch kürzere, phonologisch weniger komplexe Ausdrücke äußert: Waren früher längere Intensivierer wie \textit{fürchterlich}, \textit{schrecklich} oder \textit{wahnsinnig} gebräuchlich, so sind die neueren (auch hier zum Teil entlehnten) Intensivierer \textit{krass}, \textit{mega} oder \textit{übelst} morphologisch erheblich kürzer und von der Lautstruktur her wesentlich einfacher. Längere Intensivierer werden meinem subjektiven Empfinden nach unterdessen gänzlich vermieden, was ebenso auf einen Wandel in der internen Gestalt von Intensitätspartikeln hindeutet. Diese Ideen können vereinzelt oder auch in Kombination als Grundlage für weitere wissenschaftliche Untersuchungen dienen und bieten somit wertvolle Impulse für die zukünftige Forschung zur \mbox{sprachlichen Intensität}.


%Da Sprache im Allgemeinen und lexikalische Intensivierer im Besonderen fortwährendem Wandel ausgesetzt sind, sind progressive Untersuchungen umso bedeutender und sollten in der empirischen Sprachwissenschaft keineswegs übergangen werden. Auch wenn sich die Dynamizität von Intensitätspartikeln weder unterbinden noch beeinflussen lässt, kann sie und der durch sie hervorgerufene Status quo als Spiegel des Sprachgebrauchs einer Gesellschaft dienen und wertvolle Einblicke in die Hintergründe und Komplexität von Sprache und Kommunikation geben. 











%Um die Arbeit mit den Worten Buddhas (XXX) zu beenden: 
%
%\begin{center}
%\glqq{}In steter Veränderung ist diese Welt. Wachstum und Zerfall ist ihre wahre Natur. Die Dinge erscheinen und lösen sich wieder auf. Glücklich, wer sie friedvoll einfach nur betrachtet.\grqq{} \\
%()
%\end{center}


%Um die Arbeit mit den Worten Buddhas (XXX) zu beenden: \glqq{}In steter Veränderung ist diese Welt. Wachstum und Zerfall ist ihre wahre Natur. Die Dinge erscheinen und lösen sich wieder auf. Glücklich, wer sie friedvoll einfach nur betrachtet\grqq{}. 



%Dies hat zur Folge, dass sie in formellen Situationen und Textsorten weniger angemessen und somit vor allem in der Schriftsprache weniger gebräuchlich sind und in erster Linie auf die Spontansprache beschränkt zu sein scheinen, wie \citet[3]{Hauschild.1899} folgendermaßen festhält: \glqq{}Ihre frische, lebendige Anschaulichkeit steht in wohlthuendem Gegensatze zu der saft- und kraftlosen Art der Verstärkungswörter unserer Schriftsprache\grqq{}. 







